\section{Einleitung}
In dieser Arbeit wird die Effektivität von Retrieval-Augmented Generation (RAG) \parencite{Lewis.2020} untersucht.
\subsection{Problemstellung}
% Wissensintensive Aufgaben in Unternehmen und die Grenzen generischer LLMs (Halluzinationen, Knowledge Cut-off, fehlende Domänenexpertise)
Große Sprachmodelle (engl.\ \textit{Large Language Models}, LLMs) haben sich innerhalb weniger 
Jahre von experimentellen Forschungsartefakten zu einer technologischen Plattform entwickelt, 
die das Fundament für eine neue Generation von Anwendungen und Diensten bildet \autocite[72]{OLeary.2024}. 
Insbesondere im Bereich des betrieblichen Wissensmanagements (\textit{enterprise knowledge management}) 
fungieren sie als Katalysator für eine strukturelle Neuausrichtung, indem sie interne 
Wissensbestände über natürliche Sprache zugänglich machen und Wissensarbeit unterstützen 
\autocite[72]{OLeary.2024}. Damit verbunden ist die Erwartung, organisationsinternes Wissen, das bislang in heterogenen 
Dokumenten, Wikis, Tickets und Datenbanken verteilt ist, über natürliche Sprache zugänglich 
und nutzbar zu machen. 
Empirische Erhebungen zeigen jedoch, dass sich ein Großteil 
der Unternehmensimplementierungen weiterhin in einer experimentellen Phase befindet 
und ein signifikanter Abstand zwischen akademischen Prototypen und produktiven 
Unternehmenssystemen besteht (\textit{lab-to-market gap}) \autocite[1]{Karakurt.2025}.

Diese Erwartung steht jedoch in Spannung zu den strukturellen Eigenschaften vortrainierter 
LLMs. Erstens leiden Sprachmodelle systematisch unter sogenannten Halluzinationen, 
die im Bereich des \textit{Natural Language Processing} definiert werden als ein Phänomen,
 ,,in which the generated content appears nonsensical or unfaithful to the provided 
 source content'' \autocite[1:5]{Huang.2024}. 
 Diese Eigenschaft ist nicht nur ein praktisches Qualitätsproblem, sondern auch ein 
 theoretisches: \textcite[1]{Xu.2025} zeigen formal, dass Halluzinationen eine 
 \textit{innate limitation} von LLMs darstellen und ,,it is impossible to eliminate 
 hallucination in LLMs'' -- unabhängig von Modellarchitektur, Trainingsverfahren oder 
 Prompting-Techniken. Für Unternehmen sind unvollständige oder fabrizierte Antworten ein 
 erhebliches Risiko, da sie unmittelbar Entscheidungen beeinflussen und Schaden 
 verursachen können \autocite[1:2]{Huang.2024}.

Zweitens verfügen LLMs lediglich über parametrisches Wissen, das durch den Zeitpunkt ihres Trainings begrenzt ist (\textit{knowledge cut-off}). \textcite[1]{Cheng.2024} zeigen empirisch, dass der von Modellanbietern angegebene Cut-off nur an der Oberfläche greift und der tatsächliche, sogenannte \textit{effective cutoff} zwischen einzelnen Wissensressourcen abweichen kann. Die Autoren belegen zudem, dass selbst bei aktuelleren Trainingsdaten ein erheblicher Anteil des Modellwissens auf ältere Versionen zurückgeht: In einem untersuchten Fall stammten über 80\,\% der aus Wikipedia gelernten Dokumente aus Versionen vor 2023 \autocite[9]{Cheng.2024}. Für Unternehmen, deren interne Wissensbestände kontinuierlich evolvieren -- etwa durch neue Produkte, geänderte Prozesse oder regulatorische Anpassungen -- ist diese statische Wissensbasis unzureichend.

Drittens enthalten allgemein vortrainierte LLMs definitionsgemäß kein proprietäres oder organisationsspezifisches Wissen. Interne Dokumentationen, Kundeninformationen, technische Spezifikationen oder Compliance-Richtlinien sind nicht Teil der öffentlichen Trainingsdaten und müssen daher über gesonderte Verfahren in den Antwortprozess integriert werden \autocite[1]{Balaguer.2024}. Hinzu kommt, dass organisationsinternes Wissen in der Regel besonderen Anforderungen an Vertraulichkeit, Datenschutz und Auditierbarkeit unterliegt, die einer unkontrollierten Übertragung an externe Modellanbieter entgegenstehen.

Zur Adressierung dieser Defizite haben sich in Forschung und Praxis zwei dominante Architekturansätze etabliert. \textit{Retrieval-Augmented Generation} (RAG) verbindet ein parametrisches Sprachmodell mit einem nicht-parametrischen, externen Wissensspeicher, dessen relevante Inhalte zur Antwortzeit dynamisch abgerufen und in den Eingabekontext eingespeist werden \autocite[2]{Lewis.2020}; das Retrieval kombiniert dabei ,,LLMs' intrinsic knowledge with the vast, dynamic repositories of external databases'' \autocite[1]{Gao.2024}. \textit{Fine-Tuning}, insbesondere in parametereffizienten Varianten wie LoRA, das ,,freezes the pre-trained model weights and injects trainable rank decomposition matrices into each layer of the Transformer architecture'' \autocite[1]{Hu.2022}, integriert zusätzliches Wissen oder spezifische Verhaltensmuster hingegen direkt in die Modellparameter \autocite{Song.2025}. Wie \textcite[1]{Balaguer.2024} zusammenfassen: ,,RAG augments the prompt with the external data, while fine-Tuning incorporates the additional knowledge into the model itself''. Beide Ansätze verfolgen das Ziel der Wissensanreicherung, unterscheiden sich jedoch grundlegend in ihren technischen, organisatorischen und ökonomischen Implikationen \autocite[29]{Balaguer.2024}.

Trotz einer wachsenden Zahl vergleichender Studien fehlt in der Literatur eine konsolidierte, kriteriengeleitete Entscheidungsgrundlage, die es Praktikerinnen und Praktikern erlaubt, für einen konkreten Anwendungsfall des internen Wissensmanagements systematisch zwischen RAG, Fine-Tuning oder hybriden Ansätzen zu wählen. Bestehende Vergleiche fokussieren häufig auf einzelne Domänen \autocite[1]{Balaguer.2024,Soudani.2024} oder auf einzelne Dimensionen wie die reine Antwortgenauigkeit \autocite[2]{Ovadia.2024} und vernachlässigen organisatorische, regulatorische und wirtschaftliche Kriterien. Gleichzeitig zeigt die Enterprise-Forschung, dass viele Implementierungen sich noch in einem experimentellen Stadium befinden und der Übergang von Prototypen zu produktiven Systemen an einer Vielzahl technischer wie organisatorischer Hürden scheitert \autocite[1]{Karakurt.2025}.

Aus dieser Lücke ergibt sich der Bedarf nach einem konzeptionellen Kriterienkatalog, der die relevanten Bewertungsdimensionen aus der wissenschaftlichen Literatur ableitet und in eine praxistaugliche Entscheidungsmatrix überführt. Eine solche Matrix soll keine universelle Antwort liefern, sondern den Entscheidungsprozess strukturieren und damit dazu beitragen, dass Unternehmen ihre internen Wissensressourcen architektonisch fundiert und ressourcenschonend für den Einsatz mit LLMs nutzbar machen können.

\subsection{Zielsetzung der Arbeit}
% Entwicklung eines konzeptionellen Kriterienkatalogs

\subsection{Forschungsfrage(n)}
% F1: Welche Kriterien sind bei der Architekturauswahl zwischen RAG und Fine-Tuning für internes Unternehmenswissen relevant?
% F2: Wie lassen sich diese Kriterien in eine praxistaugliche Entscheidungsmatrix überführen?

\subsection{Methodisches Vorgehen}
% Systematische Literaturanalyse -> Kriterienableitung -> Synthese in Matrix

\subsection{Aufbau der Arbeit}